\chapter{Exact reduction contracts and the information lost by tropicalization}
\label{ch:reduction-contracts}

\reviewstatus{This chapter adds exact factorization, discrete trajectory, error,
and signed polynomial results. The companion module
\texttt{ReductionContracts.lean} contains sixteen kernel-checked declarations
with a per-declaration axiom manifest. The cited analytic-forcing obstruction
is an external theorem, not a result formalized by this module.}

Information reduction becomes a mathematical claim when the retained data
determine the quantity being predicted. This chapter supplies a necessary and
sufficient condition, then a quantitative obstruction when the condition fails.
It also shows precisely why the orders of two signed terms do not determine
the order of their sum. These statements make the project's graph, order, and
tropical interfaces testable without identifying them with a continuum proof.

\section{The exact closure criterion}

Let $X$ be a state space, $Y$ the space of retained observations, and $Z$ an
output space. A map $P:X\to Y$ describes what the representation retains;
$Q:X\to Z$ describes what a model must predict. No topology or probability
distribution is needed for the following statement.

\begin{theorem}[Factorization through retained information]
\label{thm:reduction-fiber-factorization}
Suppose $P$ is surjective. There is a map $G:Y\to Z$ satisfying
$G(Px)=Qx$ for every $x$ if and only if
\[
 Px=Py\quad\Longrightarrow\quad Qx=Qy
 \qquad(x,y\in X).
\]
When it exists, $G$ is unique.
\end{theorem}
\begin{proof}
If $G$ exists, equal observations have equal images under $G$, proving the
condition. Conversely, for each $b\in Y$ choose $x_b$ with $P(x_b)=b$ and
define $G(b)=Q(x_b)$. The fiber condition makes this independent of the
choice and gives $G(Px)=Qx$. If $G,H$ both factor $Q$, surjectivity gives,
for every $b=Px$, the equality $G(b)=Qx=H(b)$.
\end{proof}

Surjectivity is a modeling convention, not a demand that a compressed
representation cover an arbitrary ambient tensor space: one can take $Y$ to
be the image of $P$. The theorem proves existence of a set-theoretic map.
It does not prove that $G$ is continuous, measurable, Lipschitz, efficiently
computable, or representable by a particular neural architecture. Those are
separate requirements on a usable model.

For a discrete update $F:X\to X$, take $Q=P\circ F$. Then the criterion
is exactly preservation of the equivalence relation $x\sim y$ defined by
$Px=Py$. This gives a precise meaning to an exact reduced dynamics.

\begin{proposition}[Propagation of an exact discrete reduction]
\label{prop:reduction-trajectory}
Suppose $P\circ F=G\circ P$, $u_{n+1}=F(u_n)$,
$v_{n+1}=G(v_n)$, and $Pu_0=v_0$. Then $Pu_n=v_n$ for every natural
number $n$.
\end{proposition}
\begin{proof}
The initial equality is given. If it holds at $n$, then
$Pu_{n+1}=PF(u_n)=G(Pu_n)=G(v_n)=v_{n+1}$. Induction proves the claim.
\end{proof}

For a differentiable observation of a continuous-time trajectory
$\dot x=F(x)$, the chain rule instead gives
\[
 \frac{d}{dt}P(x(t))=DP(x(t))F(x(t)).
\]
The quantity to factor is therefore $Q(x)=DP(x)F(x)$. Only for a linear
observation does this become $Q(x)=PF(x)$. A nonlinear tokenizer or learned
coordinate map cannot silently use the linear rule. The factorization
theorem applies to this $Q$ as a map, while differentiability and existence
of trajectories remain analytic premises. The Lean trajectory declaration
in this chapter concerns discrete time only.

For example, $P(x,y)=x$ discards the variable needed to predict $Q(x,y)=y$.
States $(0,0)$ and $(0,1)$ share their observation and have distinct outputs.
No deterministic function of the observation alone can give both outputs.
Adding a memory state or retaining $y$ can change the representation and
remove this obstruction; additional network capacity on the same input cannot.

\section{A lower bound on every approximate closure}

The preceding obstruction has a quantitative version, independent of how
$G$ is trained or parameterized.

\begin{proposition}[Unresolved diameter bounds approximation error]
\label{prop:reduction-error-diameter}
Let $Z$ be a seminormed additive commutative group. For any
$P:X\to Y$, $Q:X\to Z$, and $G:Y\to Z$, if $Px=Py$, then
\[
 \|Qx-Qy\|\leq\|Qx-G(Px)\|+\|Qy-G(Py)\|.
\]
Consequently, if each error on the right is at most $\varepsilon$,
then $\|Qx-Qy\|\leq2\varepsilon$.
\end{proposition}
\begin{proof}
Put $g=G(Px)=G(Py)$ and write
$Qx-Qy=(Qx-g)+(g-Qy)$. The triangle inequality and
$\|g-Qy\|=\|Qy-g\|$ give the first assertion. Bounding the two terms
separately gives the second.
\end{proof}

This is a worst-case bound on a specified pair. A probabilistic average-loss
claim additionally needs a distribution and a loss; a sampled dataset does
not establish a diameter bound on every unseen fiber. The theorem provides
a diagnostic for proposed reductions: find equal retained states with
separated target stresses or derivatives. A quantum expectation readout,
graph embedding, or toric feature map is subject to the same criterion once
its accessible classical observation has been specified.

\section{Signed cancellation is invisible to orders alone}

For $p\in\mathbb Q[X]$, let $\operatorname{ord}_0(p)$ be its smallest
exponent with nonzero coefficient, with $\operatorname{ord}_0(0)=\infty$.
The finite-order statements below use nonzero inputs and nonzero residuals,
so they do not exploit a convention at zero. Tropical order calculations
are useful bookkeeping, but their additive rule must retain the possibility
of cancellation; compare the positive-semiring setting of
\citet{ReductionLitvinov2005}.

\begin{proposition}[Arbitrarily large cancellation gain]
\label{prop:reduction-signed-order}
For every integer $k\geq0$, put
\[
 p=X,\qquad q_k=-X+X^{k+2}.
\]
Then $p,q_k,p+q_k$ are nonzero and
\[
 \operatorname{ord}_0(p)=\operatorname{ord}_0(q_k)=1,
 \qquad \operatorname{ord}_0(p+q_k)=k+2>1.
\]
In particular, the identity
$\operatorname{ord}_0(p+q)=\min\{\operatorname{ord}_0p,
\operatorname{ord}_0q\}$ is false in general even for nonzero summands
and a nonzero sum.
\end{proposition}
\begin{proof}
The constant coefficients vanish. The coefficient of $X$ in $p$ is one
and that in $q_k$ is minus one, since $k+2\geq2$. Each summand has order
one and is nonzero. Their sum is exactly $X^{k+2}$, which is nonzero
and has order $k+2$. This proves all assertions.
\end{proof}

The integer $k$ is arbitrary. Thus no fixed correction to the minimum of
the input orders can recover all cancellation gains. Signed leading
coefficients must first be compared; when they cancel, higher coefficients
must be examined. For complex amplitudes, phases carry analogous information.
This issue is especially consequential for a momentum residual: its
transport, pressure, temporal, and viscous terms carry signs, and smallness
of their exact sum is a different property from the size of each term.
An order graph can organize candidate balances. A proof must still establish
the coefficient identities and bounds on the remaining terms.

\section{A current analytic obstruction to removing the force}

Theorem~1.1 of \citet{ReductionCIV2026}, version~2, addresses suitable weak
solutions smooth before the candidate singular time. It assumes a uniform
spatial $C^2$ bound on the force, spatial analyticity locally uniform on
interior time intervals, the anisotropic angular-mean bounds in its
equations~(1.7a)--(1.7b) through derivative order two, and exact axisymmetry
on a positive-radius core at every earlier time. Under these hypotheses the
candidate point is regular. Zero forcing meets the force hypotheses.
Consequently, preserving this structural class while deleting the force
cannot retain the released singularity. This is an obstruction to that
specified class, not a global unforced regularity theorem and not a
disproof of the smooth forced construction. The cited analytic theorem is
used as an external result; this extension has not formalized its proof.

For the project's future developments, this gives three concrete obligations.
A proposed information reduction must state its fibers and the outputs they
determine. A tropical residual certificate must retain signed coefficients
where leading orders tie. A transformed fluid theorem must preserve its
forcing and regularity hypotheses, or separately prove the replacements.
These obligations are compatible with graph-based proof search and quantum
learning, but they cannot be discharged by an analogy between representations.

\section{Formal coverage and reproducibility}

The companion module proves the two directions and uniqueness in
Theorem~\ref{thm:reduction-fiber-factorization}, its discrete-closure
specialization and a mismatched-fiber obstruction,
Proposition~\ref{prop:reduction-trajectory}, both inequalities in
Proposition~\ref{prop:reduction-error-diameter}, and every polynomial
assertion in Proposition~\ref{prop:reduction-signed-order}.
The declarations use ordinary set-theoretic functions, actual seminorms,
and actual rational polynomials; no placeholder predicate stands in for
the missing analytic theorem. The verification manifest records the source
hash, exact compiler and mathlib pins, process exit code, and each
declaration's foundational dependencies. Its sixteen checked declarations
are additional to the earlier finite fluid and quantum modules.
